\documentclass[10pt,letterpaper,final]{report}
\usepackage[margin=1in]{geometry}
\usepackage[utf8]{inputenc}
\usepackage{amsmath}
\usepackage{amsfonts}
\usepackage{amssymb}
\usepackage{amsthm}
\renewcommand\qedsymbol{$\blacksquare$}
\usepackage{enumerate}
\usepackage{hyperref}
\usepackage{pdfpages}
\usepackage{graphics}
\usepackage{graphicx}
\usepackage{tikz}
\usepackage{tikz-qtree}
\usepackage{forest}
\usetikzlibrary{automata,arrows}
\usepackage{mdframed}
\usepackage[
  backend=biber,
  style=numeric-comp,
  sorting=none
]{biblatex}
\addbibresource{references.bib}

\usepackage{minted}
\usemintedstyle{algol_nu}
\usepackage{xltabular}
\usepackage{pdflscape}
\usepackage{multicol}
\usepackage{underscore}
\usepackage{changepage}
\usepackage{todonotes}

\NewDocumentCommand{\sourcefile}{m v}{
  \subsection{\texttt{#2}}\label{sec:#2}
  \inputminted[linenos,breaklines,breakanywhere=true,fontsize={\fontsize{6pt}{7pt}\selectfont}]{#1}{#2}
}

% Based on template from COSC-417 at Towson University, originally authored by Marius Zimand

\renewcommand{\thesection}{\arabic{section}}

\begin{document}
\begin{center}
  \fbox{
    \vbox{
      \begin{flushleft}
        Jonathan Llovet \\  % authors' names
        EN.605.202.81 \\  %class
        2026-07-28 \\  % date
        Lab 3 \\ % ID
      \end{flushleft}
      \center{\Large{\textbf{Recursion - Converting Prefix to Postfix Expressions}}}
      %\end{mdframed}
    } % end vbox
  } % end fbox
  \vline
\end{center}

\pagebreak

\noindent\textbf{Prompt:}

\noindent Using the frequency table shown below, build a Huffman Encoding Tree. Resolve ties by giving single letter groups precedence (put to the left) over multiple letter groups, then alphabetically. Do not worry about punctuation or capitalization.

\noindent For this lab only, you may use string methods to work with the groups of letters.

\noindent This program will need to be able to read three different files: a file containing the frequency table, a file containing clear text to be encoded and a file containing encoded strings.

\noindent Print out the tree by doing a preorder traversal. Print the resulting encoding generated by the tree.

\noindent An example is given below for a 3-letter alphabet. You may use any reasonable format to print the encoding.

\begin{figure}[H]
  \centering
  \begin{forest}
    for tree={draw, minimum size=2em, inner sep=5pt, l sep=20pt}
      [\texttt{XYZ:6}
        [\texttt{X:3}]
        [\texttt{YZ:3}
          [Y:1]
          [Z:1]
        ]
      ]
  \end{forest}
  \caption{Huffman Encoding Tree Example}\label{fig:abstract-syntax-tree}
\end{figure}

\begin{verbatim}
  X: 3
  Y: 1
  Z: 2
\end{verbatim}

The tree in preorder is: \texttt{XYZ: 6, X: 3, YZ: 3, Y: 1, Z: 2}

\medskip
The code is \texttt{X = 0; Y = 10; Z = 11;}
\medskip

In your Analysis, consider whether you achieved any useful data compression with this method. Compare to conventional encoding. How would your results be affected by using a different scheme to break ties, for example, if you had given precedence to alphabetical ordering and then to the number of letters in the key? What other structures did you use and why?

Be sure to review the Programming Assignment Guidelines for specific requirements for the Analysis and before submitting this assignment.
Encode the strings in the supplied ClearText.txt file, plus several others of your choice:
Decode the strings the supplied Encoded.txt file, plus some others of your choice:
Use the frequency table in the supplied FreqTable.txt file. Consider experimenting with other
frequency tables.

As a check, here is a simple string in clear text and encoded form.

\texttt{Hello World}
\texttt{1101101000010001111100011111101000000101100}

\tableofcontents

\pagebreak

\section{The Problem to Solve}

The program should satisfy the following:

\begin{enumerate}
  \item Read frequency table from input file
  \item Read encoded strings from an input file and decode them
  \item Read plaintext strings from an input file and encode them
  \item Write the results of the previous steps to a file
  \item Echo the input
  \item Print errors to the user in a readable format
  \item Write logs to a file
  \item Present usage information to the user
  \item Have no other side effects
\end{enumerate}

\section{Strategy}\label{sec:strategy}

The structure of this lab is fairly straightforward. When called as a module, \texttt{__main__.py} calls into \texttt{lab.py}, which in turn makes use of the \texttt{encode} and \texttt{decode} packages. The \texttt{encode} package makes use of a normalize package for cleaning whitespace and punctuation from input strings, and it returns the strings fully capitalized. Both \texttt{encode} and \texttt{decode} use the \texttt{tree} package, where there are a \texttt{Tree} class and a \texttt{MinHeap} class.

Program execution proceeds as follows:

\begin{enumerate}
  \item Read arguments from command line
  \item Initialize logging
  \item Open and read file contents
  \item Process the files:
        \begin{enumerate}
          \item Build a Huffman Encoding Tree from the input frequency table
          \item Extract the codes from the tree to a dictionary for easy lookup
          \item Echo the input plaintext
          \item Compress each of the lines from the plaintext file by looking up each character from the plaintext in the codes dictionary in turn
          \item Write the compressed versions to the encoded output file
          \item Echo the input ciphertext
          \item Decompress each of the lines from the encoded file by reading each bit of the line to inform traversing through Huffman Encoding tree until we reach leaf nodes, then reading the character into a buffer and returning to the root of the tree, repeating until the line is exhausted
          \item Write the decompressed versions to the decoded output file
          \item Print errors
        \end{enumerate}
\end{enumerate}

\begin{figure}[ht]
  \centering
  \includegraphics[width=0.9\linewidth]{packages.pdf}
  \caption{Package Dependencies}\label{fig:package-dependencies}
\end{figure}


\clearpage\pagebreak

\section{Structure}\label{sec:structure}

\begin{verbatim}
lab3
|-- __init__.py
|-- __main__.py
|-- huffman
|   |-- __init__.py
|   |-- decode.py
|   |-- encode.py
|   |-- normalize.py
|   `-- tree.py
|-- lab3.py
|-- launch.json
`-- tests
    `-- huffman
        |-- __init__.py
        |-- test_decode.py
        |-- test_encode.py
        |-- test_normalize.py
        `-- test_tree.py

4 directories, 14 files

\end{verbatim}

\section{Design Decisions}

\begin{minted}[linenos,breaklines,breakanywhere=true,fontsize={\fontsize{7pt}{8pt}\selectfont}]{python}
\end{minted}

A few notable decisions: I used two different tree implementation strategies here. In the \texttt{Tree} class, the implementation uses a linked structure, because I wanted to have easy extensibility and navigation through the tree, and because I wanted to be able to repurpose the generator patterns I learned in the last lab to very simply implement the tree traversal. This ended up being a core data structure I used for my Huffman Encoding Tree, and it was used in direct conjunction with the priority queue that I implemented with a MinHeap that used an array implementation. There, I wanted to use the array implementation to simplify the enqueue and dequeue operations, which require random access to perform efficiently. I had considered using a linked list for the priority queue, but I decided that the implementation would be too slow, because of the linear time required to traverse the linked list. With the MinHeap implementation of the priority queue, the queue and enqueue operations have $O(\log n)$ time, rather than linear time, which is a huge improvement!

\section{Algorithms, Compression}

While the string encoding of the compressed data does not appear to be smaller than the plaintext data provided as input, this is really a misleading consequence of the presentation and that we're outputting our compressed values themselves as strings. If we were to store the compressed versions in binary directly and compare the length of the binary data of the plaintext and the ciphertext, we could see more plainly that the compression results in a large amount of saving. Without getting into the minutiae of typical encodings, it's worthwhile to say that Unicode encodings are represented usually with multi-bye encodings, and with encoding schemes like UTF-8, can have variable-length encodings. \cite{JoelonSoftware2003AbsoluteMinimumEvery} For ASCII characters that are representable within the space of one byte, the compression we achieve here is at most of five bits, but if the frequency table we used included characters from other alphabets that have longer representations in the two to six byte range, the compression could be much more pronounced. If a four-byte character were compressed to three bits, then that represents a saving twenty-one bits per instance of that character. A question that I'm wondering about and would like to research some more is how this kind of compression compares to mathematical representations and ideas of compressibility such as Kalmogorov complexity. I suspect that there might be other representations that could be used that would result in more compact versions of the same data, but they might use other techniques. A related set of questions arise for me out of interest in cryptography and cybersecurity. Because this encoding is based on a frequency table that can be computed from a corpus or arbitrarily chosen document, this could be used to encrypt data. I wonder how this compares to encryption methods used in industry otherwise and whether this could be subject to successful attacks based on frequency analysis or some other method. Prima facie, I am skeptical that it could be.

On the question of alternative methods for tie breaking: I don't think that ordering alphabetically first and then by score/frequency would work. The indexing is done in principle based on the frequency of the occurrence of the characters in the reference document that is the source for the frequency table, and the effect of this is to produce a method for generating short codes for characters that appear frequently, longer codes for ones that occur infrequently, so that the compression can save as much space as possible. While it might be possible to get \emph{some} encoding, I suspect that the tie breaking would not actually work, nor make sense, and that furthermore the encoding would not produce an effective compression.

\pagebreak

\section{Testing}\label{sec:testing}

As in my previous labs, I used test-driven development to build the program overall. The place where this gave me the most benefit was in building the foundational data structures for the program. I was not able to implement tests as completely as I had in previous labs, but they still provided me a much needed safety net while refactoring and iterating on my design. Over the course of the lab, I stripped out several of the layers of package design, once I realized that they were overcomplicating the code. Having tests here, as before, helped me make those structural changes.

My workflow is approximately described below. In practice, I occasionally deviated from this workflow, and at times this caused me issues. For the tree package I used this workflow in a purer form. There the tests were easier to formulate, and the bugs were (relatively) straightforward to identify. At higher levels in the stack, I was less disciplined about this, to my chagrin.

\noindent The idealized workflow for adding features is:

\begin{enumerate}
  \item Write a test that reflects an atomic aspect of the expected behavior.
  \item Validate that the test fails without the implementation in place.
  \item Write a minimal implementation that makes the test pass.
  \item Iterate with new tests.
\end{enumerate}

\noindent The complementary idealized workflow for making changes, including deletions is:

\begin{enumerate}
  \item Run all the tests and validate they all pass.
  \item Make a targeted change.
  \item Re-run all the tests and validate they all pass.
  \item If a test fails, fix the issue while preserving the desired change.
  \item Iterate.
\end{enumerate}

\noindent The third kind of workflow, for changing specifications, is:

\begin{enumerate}
  \item Run all the tests and validate they all pass.
  \item Make a targeted change to the specification.
  \item Re-run all the tests.
  \item If a test fails, fix the issue while preserving the desired change to the specification.
  \item Iterate.
\end{enumerate}

{\small
  \begin{verbatim}
python -m unittest discover -s lab3/tests
...........................
----------------------------------------------------------------------
Ran 27 tests in 0.001s

OK
\end{verbatim}
}

\pagebreak
\section{Demo}\label{sec:demo}

For the following, I assume that I start in the root of the Github repository. From some appropriate location on your filesystem, clone the repository and move into the root of the repo.
  {\small
    \begin{verbatim}
  git clone https://github.com/jllovet/jhu-en-605-202-su26-data-structures.git
  cd jhu-en-605-202-su26-data-structures
\end{verbatim}
  }

\noindent From there, move into the folder for this lab.

  {\small
    \begin{verbatim}
  cd lab-3
\end{verbatim}
  }

\noindent From here, the python module can be run as follows.

{\small
  \begin{verbatim}
  python -m lab3 \
  --level DEBUG \
  --logfile demo.lab3.log \
  --frequency_table_file ./resources/input/FreqTable.txt \
  --plaintext_file ./resources/input/ClearText.txt \
  --compression_results_file ./resources/output/EncodingResults.txt \
  --compressed_file ./resources/input/Encoded.txt \
  --decompression_results_file ./resources/output/DecodingResults.txt
\end{verbatim}
}

\noindent If the program is run with the incorrect number of arguments, then an error message similar to the one shown below will be printed to the user.

\begin{minted}[breaklines, breakanywhere, fontsize=\footnotesize]{text}
usage: python3.14 -m lab3 [-h] [--frequency_table_file FREQUENCY_TABLE_FILE] [--plaintext_file PLAINTEXT_FILE]
                          [--compression_results_file COMPRESSION_RESULTS_FILE] [--compressed_file COMPRESSED_FILE]
                          [--decompression_results_file DECOMPRESSION_RESULTS_FILE] [-l {INFO,WARNING,ERROR,DEBUG,CRITICAL,FATAL}] [-f LOGFILE]

Use Huffman Encoding to compress and decompress data

options:
  -h, --help            show this help message and exit
  --frequency_table_file FREQUENCY_TABLE_FILE
                        Pathname for the file containing the mapping from characters to frequencies
  --plaintext_file PLAINTEXT_FILE
                        Pathname for the plain text that is to be compressed
  --compression_results_file COMPRESSION_RESULTS_FILE
                        Pathname for the filename for the ciphertext resulting from compression
  --compressed_file COMPRESSED_FILE
                        Pathname for input ciphertext that is to be decompressed
  --decompression_results_file DECOMPRESSION_RESULTS_FILE
                        Pathname for the filename for the plaintext resulting from decompression
  -l, --level {INFO,WARNING,ERROR,DEBUG,CRITICAL,FATAL}
                        Sets the level of the logger. Default INFO
  -f, --logfile LOGFILE
                        Sets the filename where logs are written
\end{minted}


\noindent If the program is run with a frequency table that has invalid lines, then an error message similar to the one shown below will be printed to the user.

\begin{minted}[breaklines, breakanywhere, fontsize=\footnotesize]{text}
WARNING: 1 errors found during run!
--------------------------------------------------------------------------------
ERROR: The frequency table has invalid entries. Validate that the entries are of the form 'A - 1'

\end{minted}


\noindent A demo of multiple successive runs are shown below:

\inputminted[breaklines, breakanywhere, fontsize=\footnotesize, frame=single, fontsize=\footnotesize]{text}{resources/demo.txt}

\begin{landscape}
\begin{center}
  \begin{minipage}[t]{0.6\textwidth}
    \centering
    \textbf{\texttt{resources/input/ClearText.txt}}\\
    \inputminted[linenos, breaklines, breakanywhere, frame=single,fontsize=\footnotesize]{text}{resources/input/ClearText.txt}
  \end{minipage}\hfill
  \begin{minipage}[t]{0.6\textwidth}
    \centering
    \textbf{\texttt{resources/output/EncodingResults.txt}}\\
    \inputminted[linenos, breaklines, breakanywhere, fontsize=\footnotesize, frame=single, fontsize=\footnotesize]{text}{resources/output/EncodingResults.txt}
  \end{minipage}
\end{center}

Here the encoded data from the previous run is shown decoded. Note that the text is still in a normalized form.

\begin{center}
  \begin{minipage}[t]{0.6\textwidth}
    \centering
    \textbf{\texttt{resources/input/Encoded.txt}}\\
    \inputminted[linenos, breaklines, breakanywhere, frame=single,fontsize=\footnotesize]{text}{resources/input/Encoded.txt}
  \end{minipage}\hfill
  \begin{minipage}[t]{0.6\textwidth}
    \centering
    \textbf{\texttt{resources/output/DecodingResults.txt}}\\
    \inputminted[linenos, breaklines, breakanywhere, fontsize=\footnotesize, frame=single, fontsize=\footnotesize]{text}{resources/output/DecodingResults.txt}
  \end{minipage}
\end{center}

\pagebreak

\begin{center}
  \begin{minipage}[t]{0.6\textwidth}
    \centering
    \textbf{\texttt{resources/input/ClearText_00.txt}}\\
    \inputminted[linenos, breaklines, breakanywhere, frame=single,fontsize=\footnotesize]{text}{resources/input/ClearText_00.txt}
  \end{minipage}\hfill
  \begin{minipage}[t]{0.6\textwidth}
    \centering
    \textbf{\texttt{resources/output/EncodingResults_00.txt}}\\
    \inputminted[linenos, breaklines, breakanywhere, fontsize=\footnotesize, frame=single, fontsize=\footnotesize]{text}{resources/output/EncodingResults_00.txt}
  \end{minipage}
\end{center}

\medskip

Here the encoded data from the previous run is shown decoded. Note that here too the text is still in a normalized form.
\begin{center}
  \begin{minipage}[t]{0.6\textwidth}
    \centering
    \textbf{\texttt{resources/output/EncodingResults_00.txt}}\\
    \inputminted[linenos, breaklines, breakanywhere, frame=single,fontsize=\footnotesize]{text}{resources/output/EncodingResults_00.txt}
  \end{minipage}\hfill
  \begin{minipage}[t]{0.6\textwidth}
    \centering
    \textbf{\texttt{resources/output/ReversedDecodingResults.txt}}\\
    \inputminted[linenos, breaklines, breakanywhere, fontsize=\footnotesize, frame=single, fontsize=\footnotesize]{text}{resources/output/ReversedDecodingResults.txt}
  \end{minipage}
\end{center}
\end{landscape}

\pagebreak
\section{Enhancements and What I Learned}\label{sec:enhancements-and-what-I-learned}

\noindent \textbf{What I Leanred: Implementation of MinHeap}

I used a \texttt{MinHeap} to build a priority queue to use in the building of the Huffman Encoding Tree. While doing so I found myself constrained by my previous decision to create an \texttt{EncodingData} type. How could I do the comparisons in the \texttt{MinHeap} between nodes if I needed to compare the \texttt{EncodingData}? I solved this by writing comparison operators in dunder methods in the \texttt{EncodingData} class. By defining the methods directly into the class, I was able to provide comparisons with tie breakers, per our requirements. Importantly, these comparison methods allow the objects of this type to be compared with the built-in comparison operators like \texttt{<}, which made the implementation of the priority queue with a \texttt{MaxHeap} much easier.

\begin{minted}[breaklines, breakanywhere, fontsize=\footnotesize]{python}
def __lt__(self, other):
    # Sorting: score has highest priority, then length, then alphabetic order
    # Score
    if self.score < other.score:
        return True
    if self.score > other.score:
        return False
    # implies that self.score == other.score
    # Length
    if len(self.characters) < len(other.characters):
        return True
    if len(self.characters) > len(other.characters):
        return False
    # implies len(self.characters) == len(other.characters)
    # Alphabetic order
    return self.characters < other.characters
\end{minted}

\noindent \textbf{Enhancement: Unit Testing}
\medskip

During this lab I learned a valuable lesson about problem specification and how it can bite you. When writing the comparison operators for the minheap, I had misread my notes and the ZyBook about the implementation, so that I was swapping a node with its largest child while percolating down, rather than with its smallest child. This caused the program to run incorrectly, but my tests were passing with a specification I had entered that was incorrect. It took until later on in my development to realize this mistake, and I had to update my tests and manually walk back through many of my test cases for dequeueing operations. The conventional wisdom is to measure twice and cut once. An additional layer to this should be that we must also think twice about what we are measuring! This was valuable to re-encounter as a potential problem to run into, because one of my research interests is in formal verification and using proof assistants. They represent in some ways tools that are similar to tests of the kind here, but with much stronger guarantees. It's all the more important, therefore, to validate that whatever models we make with them we also double check. Put otherwise, a test may validate that something does not have an error according to a given specification, but it cannot validate that the specification is the right thing in the first place. Similarly with a proof, it can validate that a solution is valid up to the model, but what it cannot do is validate that that is in fact the right thing to have modeled for a given problem or use case. That takes additional human judgement.

For more details of the unit testing that I did as an enhancement, please see \ref{sec:testing}.

\medskip
\noindent \textbf{Enhancement: Logging}
\medskip

Here, as in my previous lab, I added logging that allows one to see the program execution information at varying levels of detail. There is a flag on the CLI for the module to specify a logging level, and the user can additionally specify a logfile that they would like to write to. One of the things that I learned from this round of development is that in addition to my tests, I would have liked to have had logging built in from earlier on in the project. When I added it in, I realized that I missed the easy visibility that it provides while working on building features. Something I still want to explore, however, is the use of monads and types to have my functions be pure, rather than having out of band side effects like writing to a logger.

Some example logs from a successful run are below:

\pagebreak

\inputminted[breaklines, breakanywhere, fontsize=\footnotesize, frame=single]{text}{demo.lab3.log}

\pagebreak

\section{Resouces Referenced While Coding}

While developing this lab, I referenced the following resources:

\begin{enumerate}
  \item \textcite{zybooks-ds2026}
  \item \textcite{PythondocumentationOperatorStandard}
  \item \textcite{StackOverflowHowToUseDunderMethodsForObjectComparison}
  \item \textcite{EnwikipediaorgTreeTraversal}
  \item \textcite{PythonEnhancementProposalsPEP0636}
  \item \textcite{JoelonSoftware2003AbsoluteMinimumEvery}
\end{enumerate}

\pagebreak

\section{Source Code}
The full source code, including LaTeX, tooling, other artifacts, and git history is available on my Github repository for the class here:
\begin{center}
  \url{https://github.com/jllovet/jhu-en-605-202-su26-data-structures}
\end{center}

\printbibliography[heading=bibintoc]

\end{document}